% GENERATED by scripts/manual_tables.py from the adapter registry. Do not edit:
% rerun the script. Hand-written text: blocks/detection/signal_threshold/
\ubblock{detection.signal_threshold}%
  {Signal threshold to spans}% what the app calls it
  {Threshold (Signal -> SpanSet)}% registry display name
  {detect}% insert-modal tab
  {Signal}% input
  {SpanSet}% output
  {Adapters/detection\_signal\_threshold.py}% code
  {detection/signal_threshold}% hand-written folder under blocks/
\ubblockdescription{detection.signal_threshold}{Generic block: cuts the chain's current signal at an absolute level in its own units and keeps the contiguous runs above it, below it, or past it on either side of zero, written to `detections`. Put baseline removal before it; a span scores by how far past the cut it reaches. Optionally join spans closer than `merge\_gap\_s` and drop those shorter than `min\_duration\_s` (both off at 0).}
\ubblockflags{detection.signal_threshold}{%
}
\ubblocksideinputs{detection.signal_threshold}{%
}
\ubblockparams{detection.signal_threshold}{%
  \ubparam{threshold}{0.0}{any number}{The cut: an absolute level in the signal's own units as stored (M2 is stored in volts, so -0.005 reads as -5 mV)}%
  \ubparam{direction}{\texttt{above}}{one of \texttt{above}, \texttt{below}, \texttt{either}}{Which side of the cut counts: `above` keeps x > threshold, `below` keeps x < threshold (a drop after baseline removal), `either` keeps |x| > |threshold|}%
  \ubparam{min\_\allowbreak{}duration\_\allowbreak{}s}{0.0}{number, at least 0}{Drop a span shorter than this many seconds (0 keeps every run, however short); applied after merging}%
  \ubparam{merge\_\allowbreak{}gap\_\allowbreak{}s}{0.0}{number, at least 0}{Join spans separated by less than this many seconds into one, scoring the max of its parts (0 never joins); applied before the duration floor}%
}
